\section{Introducción}

La producción audiovisual en directo ha experimentado en los últimos años una profunda
transformación, impulsada por la democratización del hardware de bajo coste y el
crecimiento exponencial del \textit{streaming} en Internet. Eventos técnicos de referencia como
el Chaos Communication Congress (CCC), organizado por el Chaos Computer Club en Alemania,
han puesto de manifiesto la necesidad de disponer de herramientas de realización
profesional que sean al mismo tiempo económicas, fiables y auditables.

En este contexto surge Voctomix, desarrollado por el equipo C3VOC (\textit{Chaos Communication
Congress Video Operations Crew}). Frente a soluciones previas como \textit{dvswitch} ---limitada
en resolución y escalabilidad--- o herramientas comerciales como vMix o OBS Studio
---diseñadas para un único operador en un único equipo---, Voctomix adopta una arquitectura
cliente-servidor desacoplada que separa el motor de mezcla multimedia de la interfaz
de control, permitiendo distribuir la carga de procesamiento en distintos nodos de red.

El presente Trabajo de Fin de Grado toma como punto de partida la versión original
de Voctomix y propone una serie de extensiones orientadas a su uso en producciones
en directo de mediana escala: rotulación dinámica sobre el vídeo, gestión avanzada
del blanking de stream, seguimiento automático del audio a la fuente de vídeo activa,
telemetría en tiempo real mediante un servicio de notario, y la contenerización completa
del sistema mediante Docker y Docker Compose. El resultado es Voctomix 2.0, una
plataforma más versátil, desplegable y adaptada a flujos de trabajo profesionales.

% --- Figura sugerida ---
% \begin{figure}[H]
%   \centering
%   \includegraphics[width=0.85\textwidth]{figuras/gui_voctomix2.png}
%   \caption{Interfaz gráfica de Voctomix 2.0 en funcionamiento.}
%   \label{fig:gui_voctomix2}
% \end{figure}

% --- Figura sugerida ---
% \begin{figure}[H]
%   \centering
%   \includegraphics[width=0.85\textwidth]{figuras/diagrama_contexto.png}
%   \caption{Diagrama de contexto: Voctomix en una producción audiovisual típica.}
%   \label{fig:diagrama_contexto}
% \end{figure}
